\section{Model}\label{sec:model}

\subsection{Random parameters negative binomial model}\label{sec:rpnb}

Let $y_{it}$ be the number of crashes at site $i=1,\dots,n$ in period $t=1,\dots,T_i$, with $N=\sum_i T_i$ observations in total. Cross-sectional data correspond to $T_i=1$. Let $\x_{it}=(x_{it1},\dots,x_{itp})^\top$ collect $p$ candidate covariates and let $o_{it}$ be a known offset, for example the logarithm of segment length times the length of the period. Conditional on a rate $\lambda_{it}$, the count follows a negative binomial distribution with mean $\lambda_{it}$ and dispersion $r>0$,
\begin{equation}\label{eq:nb}
\Pr(y_{it}=y \mid \lambda_{it}, r)=\frac{\Gamma(y+r)}{\Gamma(r)\,y!}\left(\frac{r}{r+\lambda_{it}}\right)^{r}\left(\frac{\lambda_{it}}{r+\lambda_{it}}\right)^{y},
\end{equation}
so that $\Var(y_{it}\mid\lambda_{it})=\lambda_{it}+\lambda_{it}^2/r$.
The random parameters specification is
\begin{equation}\label{eq:rp}
\log\lambda_{it}=o_{it}+\alpha+\sum_{j=1}^{p}x_{itj}\,\beta_{ij},
\qquad \beta_{ij}=\beta_j+\sigma_j z_{ij},
\qquad z_{ij}\overset{\text{iid}}{\sim}\Normal(0,1),
\end{equation}
so that the coefficient of covariate $j$ at site $i$ is normally distributed with mean $\beta_j$ and standard deviation $|\sigma_j|$. The draw $z_{ij}$ is shared by all periods of site $i$, so in panel data a site keeps its coefficient over time, while in cross-sectional data each observation has its own coefficient. Substituting the second equation into the first gives the non-centered form
\begin{equation}\label{eq:noncentered}
\log\lambda_{it}=o_{it}+\alpha+\x_{it}^\top\bbeta+\varepsilon_{it},
\qquad \varepsilon_{it}=\sum_{j=1}^{p}\sigma_j\,x_{itj}\,z_{ij}.
\end{equation}
We allow $\sigma_j$ to take any real value. Its sign is not identified, because $(\sigma_j,z_{ij})$ and $(-\sigma_j,-z_{ij})$ give the same likelihood, and all reported quantities depend on $\sigma_j$ only through $|\sigma_j|$. The advantage of \eqref{eq:noncentered} is that, given $z_{ij}$, the parameter $\sigma_j$ enters the linear predictor linearly, exactly as $\beta_j$ does \citep{papaspiliopoulos2007general, fruhwirth2010stochastic}.

Each covariate is in one of three states,
\begin{equation}\label{eq:states}
\state_j=\begin{cases}
\mathrm{E}\ (\text{excluded}) & \text{if } \beta_j=0 \text{ and } \sigma_j=0,\\
\mathrm{F}\ (\text{fixed}) & \text{if } \beta_j\neq 0 \text{ and } \sigma_j=0,\\
\mathrm{R}\ (\text{random}) & \text{if } \sigma_j\neq 0 .
\end{cases}
\end{equation}
The specification problem is to learn $\state=(\state_1,\dots,\state_p)$. Some columns of $\x_{it}$ may be admissible only as fixed effects, for example quadratic terms introduced to protect against the confounding described in Section~\ref{sec:identification}. For those columns we set $\sigma_j=0$.

\begin{remark}[Centering changes the model]\label{rem:centering}
For a fixed parameter, replacing $x_{itj}$ by $x_{itj}-c$ only shifts the intercept. For a random parameter it does more: $\beta_{ij}(x_{itj}-c)=\beta_{ij}x_{itj}-c\beta_j-c\sigma_j z_{ij}$, and the last term is a site-specific random intercept. Centering therefore moves heterogeneity between the slope and the intercept and changes the variance function $v(\cdot)$ defined below. We scale continuous covariates by their standard deviations but do not center them, and we leave binary covariates coded $0/1$, which keeps the conventional interpretation of random parameters.
\end{remark}

\subsection{What crash data can reveal about random parameters}\label{sec:identification}

Write $\eta_{it}=o_{it}+\alpha+\x_{it}^\top\bbeta$, $D=\operatorname{diag}(\sigma_1^2,\dots,\sigma_p^2)$ and $v(\x)=\x^\top D\x=\sum_j\sigma_j^2x_j^2$. Given the covariates, $\varepsilon_{it}$ in \eqref{eq:noncentered} is a sum of independent normal terms, so $\varepsilon_{it}\sim\Normal\{0,v(\x_{it})\}$ and $\Cov(\varepsilon_{it},\varepsilon_{is})=\x_{it}^\top D\,\x_{is}$. The following results, proved in \ref{app:proofs}, describe what this implies for the distribution of crash counts.

\begin{proposition}[Cross-sectional data]\label{prop:cross}
Let $T_i=1$ and drop the index $t$. Then:
\begin{enumerate}[label=(\alph*)]
\item The likelihood contribution of site $i$ is
\[
\int \NB\{y_i \mid e^{\eta_i+u}, r\}\,\phi\{u \mid 0, v(\x_i)\}\,du,
\]
where $\phi(\cdot\mid 0,v)$ is the normal density with variance $v$. The likelihood therefore depends on $(\sigma_1,\dots,\sigma_p)$ only through the variance function $v(\cdot)$.
\item The conditional mean and variance are
\begin{equation}\label{eq:moments}
\begin{aligned}
\mu_i&=\E(y_i\mid\x_i)=\exp\{\eta_i+v(\x_i)/2\},\\
\Var(y_i\mid\x_i)&=\mu_i+\mu_i^2\left\{\left(1+\tfrac{1}{r}\right)e^{v(\x_i)}-1\right\}.
\end{aligned}
\end{equation}
\item A necessary condition for $(\sigma_1^2,\dots,\sigma_p^2)$ to be locally identified at a point with all $\sigma_j^2>0$ is that the functions $\x\mapsto x_j^2$, $j=1,\dots,p$, be linearly independent on the support of $\x$.
\end{enumerate}
\end{proposition}

Part (b) has three consequences that bear directly on specification practice.

\begin{corollary}\label{cor:confounding}
Under the conditions of Proposition~\ref{prop:cross}:
\begin{enumerate}[label=(\roman*)]
\item If $x_j$ is binary, then $x_j^2=x_j$ and the log mean contains $x_j(\beta_j+\sigma_j^2/2)$. The mean identifies only $\beta_j+\sigma_j^2/2$, and information about $\sigma_j$ comes solely from the difference in overdispersion between sites with $x_j=1$ and $x_j=0$.
\item If $x_j$ is continuous, a random parameter adds $\sigma_j^2x_j^2/2$ to the log mean, which is indistinguishable in the mean from a fixed quadratic effect with a positive coefficient. A convex effect omitted from a linear specification can therefore surface as a spurious random parameter. A concave effect cannot.
\item The first two conditional moments in \eqref{eq:moments} coincide with those of a fixed parameters NB model with log mean $\eta_i+v(\x_i)/2$ and covariate-dependent dispersion $r(\x_i)=\left[(1+1/r)e^{v(\x_i)}-1\right]^{-1}$. Separating a random parameter from varying dispersion \citep{miaou2003modeling} therefore relies on higher moments, that is, on the difference between lognormal and gamma mixing.
\end{enumerate}
\end{corollary}

Panel data change the picture because the same draw $z_{ij}$ affects all periods of a site.

\begin{proposition}[Panel data]\label{prop:panel}
For $t\neq s$,
\begin{equation}\label{eq:panelcov}
\Cov(y_{it},y_{is}\mid \x_{i1},\dots,\x_{iT_i})=\mu_{it}\,\mu_{is}\left\{\exp\left(\x_{it}^\top D\,\x_{is}\right)-1\right\},
\end{equation}
where $\mu_{it}=\exp\{\eta_{it}+v(\x_{it})/2\}$. The covariance does not depend on the dispersion $r$.
\end{proposition}

The negative binomial dispersion, fixed or covariate dependent, acts independently in each period and implies zero within-site covariance. The covariance in \eqref{eq:panelcov} is therefore information about $D$ that the dispersion cannot mimic. A site-level random intercept, which corresponds to a column $x_{it0}\equiv 1$, also induces within-site covariance, but a constant one. Random slopes are identified from how the covariance changes with the covariates, which is why the random intercept should be included among the candidates in panel applications.

Proposition~\ref{prop:panel} is derived under a correctly specified model. It shows that the within-site covariance carries information about $D$ that the dispersion cannot mimic, but a misspecified model need not use that information well, and two failures matter in practice. First, suppose the mean omits a term $g(x_{ij})$ in a covariate that is constant within a site, as geometric variables usually are. The omitted term is then the same in every period of site $i$, and the random component $\sigma_jx_{ij}z_{ij}$, whose draw is also shared across periods, can absorb it by setting $z_{ij}\approx g(x_{ij})/(\sigma_jx_{ij})$. Second, suppose the dispersion varies with a covariate but the model assumes a common $r$. By Proposition~\ref{prop:cross}(b), a random parameter on that covariate raises the marginal variance at the sites where the data are more dispersed, and the gain in fit can outweigh the cost of implying a within-site covariance that the data do not show. Section~\ref{sec:sim-bc} shows that both failures occur in panel data as well as in cross-sectional data.

\begin{remark}[Practical implications]\label{rem:practice}
Evidence for a random parameter is evidence of a covariate-dependent variance in the log rate. Genuine heterogeneity in the effect can produce it, but so can omitted convex curvature and varying dispersion, and panel data do not remove either confounder when the model is misspecified. We therefore include quadratic terms of continuous covariates as fixed-only candidates in cross-sectional and panel applications alike, so that curvature is not absorbed into $\sigma_j$. Varying dispersion calls for the analogous safeguard on the dispersion side, which is to model the dispersion with its own candidates, as in \eqref{eq:dispersion} below. Without it, a selected random parameter should be read as evidence of heterogeneity in either the effect or the dispersion.
\end{remark}

\subsection{Prior specification}\label{sec:prior}

We introduce indicators $\gamma_j=\mathbb{1}(\beta_j\neq0)$ and $\delta_j=\mathbb{1}(\sigma_j\neq0)$ and use point-mass spike-and-slab priors for both:
\begin{align}
\beta_j\mid\gamma_j,\tau_\beta &\sim \gamma_j\,\Normal(0,\tau_\beta^2)+(1-\gamma_j)\,\mathbb{1}_{\{0\}}, \label{eq:prior-beta}\\
\gamma_j\mid\pi_\beta &\sim \Bern(\pi_\beta), \nonumber\\
\sigma_j\mid\delta_j,\tau_\sigma &\sim \delta_j\,\Normal(0,\tau_\sigma^2)+(1-\delta_j)\,\mathbb{1}_{\{0\}}, \label{eq:prior-sigma}\\
\delta_j\mid\gamma_j,\pi_{\sigma1},\pi_{\sigma0} &\sim \Bern\{\gamma_j\pi_{\sigma1}+(1-\gamma_j)\pi_{\sigma0}\}. \nonumber
\end{align}
The prior for $\delta_j$ encodes heredity \citep{chipman1996bayesian}. Under \emph{strong heredity}, $\pi_{\sigma0}=0$ and a parameter can be random only if its mean is nonzero, which matches the way specifications are built in practice. Under \emph{weak heredity}, $\pi_{\sigma0}$ is small but positive, which admits effects whose mean is zero but whose sign varies across sites. The three states in \eqref{eq:states} correspond to $(\gamma_j,\delta_j)=(0,0)$, $(1,0)$ and $\delta_j=1$.

The inclusion probabilities receive beta hyperpriors, $\pi_\beta\sim\Betad(a_\beta,b_\beta)$, $\pi_{\sigma1}\sim\Betad(a_\sigma,b_\sigma)$ and, under weak heredity, $\pi_{\sigma0}\sim\Betad(a_0,b_0)$. Learning these probabilities from the data produces the multiplicity adjustment of \citet{scott2010bayes}: as more candidates are considered, the prior probability that any particular one is random falls in line with the evidence that most are not. The slab scales $\tau_\beta$ and $\tau_\sigma$ are fixed at substantive values rather than learned. When the slab scale and the inclusion probability are both learned, a posterior with many coefficients in a narrow slab and a posterior with few coefficients in a wide slab fit the data almost equally well, and in small samples the sampler drifts toward the first, including nearly every candidate. We observed exactly this behavior in pilot runs with half-Cauchy priors on the slab scales, which we therefore retain only as a sensitivity option (Section~\ref{sec:computation}). A Gaussian slab on the signed $\sigma_j$ implies $\sigma_j^2\mid\delta_j=1\sim\tau_\sigma^2\chi^2_1$, a prior on the random parameter variance that places substantial mass near zero. Vague inverse-gamma priors on variances, common in applied Bayesian crash modeling, have vanishing density at zero and tend to overstate heterogeneity when the data are weak \citep{gelman2006prior}. Finally, $\alpha\sim\Normal(0,s_\alpha^2)$ and $r\sim\Gammad(a_r,b_r)$.

\paragraph{Default hyperparameters} With covariates scaled to unit standard deviation we use $a_\beta=b_\beta=a_\sigma=b_\sigma=1$, $\tau_\beta=\tau_\sigma=0.5$, $s_\alpha=10$, $a_r=1$ and $b_r=0.1$. A slab standard deviation of $0.5$ for $\beta_j$ makes rate ratios between about $0.37$ and $2.7$ per standard deviation of a covariate plausible, which covers the effects reported in safety performance functions. With $|\sigma_j|=0.5$, the rate ratio for a one-standard-deviation increase in $x_j$ at a site one standard deviation above or below the average effect differs from the average rate ratio by a factor of $e^{\pm0.5}$, about $0.61$ to $1.65$, which is already substantial heterogeneity, so $\tau_\sigma=0.5$ makes such values typical without ruling out larger ones. \tbd{Sensitivity analysis to $(a_\sigma,b_\sigma)$, $\tau_\sigma$ and heredity in Sections~\ref{sec:simulation} and~\ref{sec:application}.}

\paragraph{Covariate-dependent dispersion} Section~\ref{sec:sim-bc} shows that a dispersion that varies with a covariate is mistaken for a random parameter on that covariate unless the dispersion itself is modeled. The full version of the method therefore lets the dispersion depend on candidate covariates $\bm w_{it}$, which in our experiments are the same as the candidates for the mean:
\begin{equation}\label{eq:dispersion}
\begin{gathered}
\log r_{it}=\rho_0+\bm w_{it}^\top\bm\rho,\\
\rho_k\mid\eta_k\sim\eta_k\,\Normal(0,\tau_\rho^2)+(1-\eta_k)\,\mathbb{1}_{\{0\}},\qquad
\eta_k\mid\pi_\rho\sim\Bern(\pi_\rho),
\end{gathered}
\end{equation}
with $\pi_\rho\sim\Betad(1,1)$, $\tau_\rho=0.5$ and $\rho_0\sim\Normal(0,2^2)$. The dispersion is bounded below, $r_{it}=\max\{\exp(\rho_0+\bm w_{it}^\top\bm\rho),\,0.01\}$, which allows a variance of up to $\mu+100\mu^2$, far beyond what crash data show, and keeps the P\'olya-Gamma draws of Section~\ref{sec:computation} well defined. Extra variation linked to a covariate then has three competing explanations within one model: a random parameter on the covariate, a dispersion effect of the covariate and, through the squared candidates, curvature in the mean. The posterior weighs them against each other, and a candidate is declared to have a dispersion effect if $\Pr(\eta_k=1\mid\y)>1/2$.

\subsection{Posterior summaries for practice}\label{sec:summaries}

For each candidate the sampler returns the posterior state probabilities
\begin{equation}\label{eq:stateprob}
\begin{aligned}
P_j(\mathrm{E})&=\Pr(\gamma_j=0,\delta_j=0\mid\y),\\
P_j(\mathrm{F})&=\Pr(\gamma_j=1,\delta_j=0\mid\y),\\
P_j(\mathrm{R})&=\Pr(\delta_j=1\mid\y).
\end{aligned}
\end{equation}
Our default decision applies the median probability rule of \citet{barbieri2004optimal} to each indicator: a candidate is declared random if $P_j(\mathrm{R})>1/2$, and otherwise fixed if $\Pr(\gamma_j=1\mid\y)>1/2$. Under a zero-one loss for each covariate the Bayes decision would instead be the state with the largest posterior probability. In the simulation study of Section~\ref{sec:simulation} the two rules disagreed for about 3\% of candidates. When the claim of interest is that a parameter is random, the set $\widehat{\mathcal R}_c=\{j:P_j(\mathrm{R})\ge c\}$ has estimated Bayesian false discovery rate $\sum_{j\in\widehat{\mathcal R}_c}\{1-P_j(\mathrm{R})\}/|\widehat{\mathcal R}_c|$, and $c$ can be chosen as the smallest threshold that keeps this rate below a target $q$ \citep{newton2004detecting}. Effect summaries are averaged over specifications: the posterior distributions of $\beta_j$ and $|\sigma_j|$, the share of sites with a positive effect $\Phi(\beta_j/|\sigma_j|)$ for parameters in state R, and site-specific coefficients $\beta_{ij}$, which are shrunk toward $\beta_j$ according to the information in each site's data.
